\section*{Capítulo \ref{ch:g2:from-seed-to-plant} --- Da semente à planta}

\begin{solution}{exo:g2:from-seed-to-plant:1}
Uma plantinha já começada (o germe da \omterm{def:g1:plants-around-us:plant}{planta}, com o início de uma
\omterm{def:g1:plants-around-us:parts}{raiz} e de um broto) e uma reserva de alimento para alimentar essa
plantinha.
\end{solution}

\begin{solution}{exo:g2:from-seed-to-plant:2}
\omterm{def:g2:from-seed-to-plant:germination}{Germinar} é acordar e começar a crescer: a \omterm{def:g2:from-seed-to-plant:seed}{semente} incha, a casca se
abre e ela começa a virar \omterm{def:g1:plants-around-us:plant}{planta}. A \omterm{def:g1:plants-around-us:parts}{raiz} sai primeiro e mergulha
para baixo.
\end{solution}

\begin{solution}{exo:g2:from-seed-to-plant:3}
Água e calor. A luz não está na lista --- a \omterm{def:g2:from-seed-to-plant:seed}{semente} \omterm{def:g2:from-seed-to-plant:germination}{germina} numa boa
no escuro, debaixo da terra.
\end{solution}

\begin{solution}{exo:g2:from-seed-to-plant:4}
Ela vive da reserva de alimento guardada dentro da \omterm{def:g2:from-seed-to-plant:seed}{semente} --- as
duas metades gordinhas do feijão. Essa reserva alimenta a plantinha
até as \omterm{def:g1:plants-around-us:parts}{folhas} verdes chegarem à luz e ela conseguir se alimentar
sozinha.
\end{solution}

\begin{solution}{exo:g2:from-seed-to-plant:5}
Por exemplo: feijão, ervilha, lentilha, arroz, \omterm{def:g2:from-seed-to-plant:seed}{semente} de girassol,
caroços de maçã, caroço de pêssego.
\end{solution}

\begin{solution}{exo:g2:from-seed-to-plant:6}
Porque a \omterm{def:g2:from-seed-to-plant:seed}{semente} precisa de calor além de água: na terra fria do
inverno ela ficaria esperando, ou apodreceria, em vez de \omterm{def:g2:from-seed-to-plant:germination}{germinar}. A
primavera traz a terra quentinha que a \omterm{def:g2:from-seed-to-plant:seed}{semente} está esperando.
\end{solution}

\begin{solution}{exo:g2:from-seed-to-plant:7}
\omterm{def:g2:from-seed-to-plant:seed}{Semente}, \omterm{prop:g2:from-seed-to-plant:cycle}{plântula}, \omterm{def:g1:plants-around-us:plant}{planta} cheia de \omterm{def:g1:plants-around-us:parts}{folhas}, \omterm{def:g1:plants-around-us:plant}{planta} com \omterm{def:g1:plants-around-us:parts}{flores}, vagem
com feijões novos.
\end{solution}

\begin{solution}{exo:g2:from-seed-to-plant:8}
As duas \omterm{def:g1:plants-around-us:parts}{raízes} crescem para baixo e os dois brotos crescem para
cima, não importa como os feijões estavam deitados. A plantinha
sempre acha o para cima e o para baixo --- exatamente o que a
observação prevê.
\end{solution}

\begin{solution}{exo:g2:from-seed-to-plant:9}
Na gaveta, as \omterm{def:g2:from-seed-to-plant:seed}{sementes} não tinham água, então não podiam \omterm{def:g2:from-seed-to-plant:germination}{germinar};
mas uma \omterm{def:g2:from-seed-to-plant:seed}{semente} boa sem água não morre --- ela espera. Um ano
depois, com água e calor, as \omterm{def:g2:from-seed-to-plant:seed}{sementes} que esperavam acordaram
normalmente.
\end{solution}

\begin{solution}{exo:g2:from-seed-to-plant:10}
A reserva de alimento da \omterm{def:g2:from-seed-to-plant:seed}{semente} --- as duas metades gordinhas do
feijão --- faz o papel da gema. O \omterm{def:g2:animals-grow-up:born}{ovo} e a \omterm{def:g2:from-seed-to-plant:seed}{semente} são os dois um
começo de um novo \omterm{def:g1:living-or-not:living}{ser vivo}, cheios do alimento de que o pequeno
precisa antes de conseguir se alimentar sozinho.
\end{solution}
